\chapter{Implementation}\label{ch:impl}

This chapter describes the implementation thematically (not as a chronological development
log): the three-repository split, the SOTA and allocator adapters, the permutation code
generation and flag system, the concept/CRTP realization of the axes, concurrency and
reclamation, the cache-coherence-aware telemetry, and the measurement pipeline.

\section{Three-Repository Architecture}\label{sec:repos}

The implementation is distributed over three Git repositories with clearly separated
responsibilities:

\begin{itemize}
  \item \texttt{comdare-cache-engine} --- \textbf{how} things are measured: the tool
        library and the builder. Key components: the builder (XML config parser,
        \texttt{generate\_module\_from\_profile} code generation, permutation loop, module
        loader, and the seven-phase \texttt{ExperimentDriver}); the ABI headers
        (\texttt{search\_engine}, \texttt{execution\_engine}, \texttt{baustein\_variants},
        \texttt{resolve\_baustein}); 30 SOTA profiles under
        \texttt{algorithm\_profiles/sota}; and the test suites. The source tree follows a
        Pitchfork-style layout (\texttt{apps/} for executables, \texttt{libs/} for
        libraries, with \texttt{libs/common} for cross-cutting code).
  \item \texttt{comdare-prt-art} --- the \textbf{probe}: the hybrid
        \texttt{PrtArtSearchEngine} (vector/map/tuple API), its ABI adapters, its building
        blocks (4+2 allocator pools, OLC with reserved blocks, multi-level layout,
        path-oriented prefetch, density tracker, hypothesis metrics H1/H2/H3), and the
        probe-side reimplementations.
  \item \texttt{probst-Diplomarbeit-cache-engine} --- \textbf{what} is tested: the
        \texttt{messung\_driver}, the experiment configurations, the evaluation pipeline,
        and the manuscript.
\end{itemize}

\section{Adapters for State-of-the-Art Algorithms and Allocators}\label{sec:adapters}

Each SOTA design and each allocator is integrated through a thin adapter that exposes the
original implementation behind the engine ABI. The adapters follow a uniform pattern: a
\texttt{COMDARE\_HAVE\_<X>} compile flag guards an \texttt{\#if defined} block that wires
the original API, with a safe fallback (e.g.\ \texttt{std::malloc} or a typed exception)
when the original is unavailable. This keeps the build self-contained while allowing the
original sources to be activated per dependency. The hierarchical-prefetch integration
(P27) additionally ships a build-time call-graph analyzer (\texttt{p27\_bundle\_finder}, a
C++23 port of the original \texttt{hp\_soft} tool) that identifies functions whose subtree
footprint exceeds a threshold as bundle entry points.

\section{Permutation Code Generation and Flag System}\label{sec:codegen}

The builder turns each profile into a pre-compiled C++23 module
(\texttt{generate\_module\_from\_profile}). Every permutation is identified by a
bit-packed \emph{permutation flag} structure. With the N-phase refinement of the matrix,
the flag system grows from 9 banks ($\sim$50 bit) to 14 banks with sub-bank bitfields and
a total \textbf{82-bit permutation identifier}, mirroring the fourteen axes and their
sub-axes. A constraint filter masks out invalid axis combinations before any module is
generated.

\section{Concept/CRTP Realization of the Axes}\label{sec:concept-crtp}

Each axis is a topic directory containing a two-layer concept (a broad topic concept and a
narrow axis concept) and a CRTP base class secured by a \texttt{requires} clause. The
probe-versus-default selection per axis is resolved at compile time via
\texttt{resolve\_baustein.hpp} (\texttt{std::conditional\_t} over tag specializations), so
the hot path contains neither \texttt{virtual} calls nor run-time switches --- the
zero-cost abstraction of Section~\ref{sec:cpp23}.

\section{Concurrency and Reclamation}\label{sec:concurrency}

Concurrency is itself an axis (8.1 pattern, 8.2 locking mode), realized through a
concurrency manager that coordinates the disciplines OLC, HTM, STM, lock-free, RCU, and
hazard pointers. Reclamation is a sub-axis of the allocator (6.2): epoch-based, RCU,
hazard-pointer, and QSBR strategies are selectable per permutation. PRT-ART contributes
OLC with reserved value blocks as its concurrency variant.

\section{Telemetry against Cache-Coherence Effects}\label{sec:telemetry}

Per-node counters cause cache-line ping-pong on multi-core systems, which distorts the
very measurements they collect. Following the Kuehn line of work (P28), the telemetry axis
therefore offers cache-coherence-aware strategies: \emph{leaf-only counters} (the main
variant), \emph{sampling} (every $N$-th leaf access), and \emph{offline recompute}
(bottom-up aggregation before reordering); the classic per-node counter is retained only
as a labeled anti-pattern for comparison.

\section{Measurement Pipeline}\label{sec:pipeline}

The evaluation tool chain is a sequence of small programs:
\texttt{sample\_data\_generator} (deterministic test data without real hardware)
$\rightarrow$ \texttt{messung\_driver} (drives the series) $\rightarrow$
\texttt{binary\_to\_csv} (deserializes the binary measurement records) $\rightarrow$
\texttt{csv\_to\_latex} (booktabs tables with per-profile algorithm fact sheets)
$\rightarrow$ \texttt{diagram\_generator} (A4-aware TikZ plots) $\rightarrow$
\texttt{latex\_to\_pdf}. The measurement path in the execution engine is compiled in only
under \texttt{-DCOMDARE\_EXPERIMENT\_MODE=ON}, so the production path carries no measurement
overhead. Continuous integration (GitLab CI plus a mirrored GitHub workflow) runs across
all three repositories.
